% ============================================================
% Delivery Note Template --- VIEWS
% ============================================================
% A delivery note accompanies a set of data handed to someone outside
% the team that produced it. It records WHAT was delivered, HOW to read
% it, and WHAT IS NOT in it. It does not ask the recipient to decide
% anything --- that is a pre-release note's job.
%
% TO USE THIS TEMPLATE:
%   1. Copy this directory to ../<delivery-id>/ and work there.
%   2. Fill the six \delivery* macros below. All six are required.
%   3. Build the \section + \input spine in the DOCUMENT BODY to suit
%      this delivery. There is no fixed structure --- comment in,
%      comment out, reorder, rename. A stub exists for each section in
%      sections/; delete the ones you do not use.
%   4. Write body prose in sections/*.tex. Those files contain NO
%      \section commands --- the headings live here, so reordering is a
%      single-file edit.
%   5. Build with ./build.sh (three pdflatex passes).
%
% Conventions: ../README.md, and the house style guides at
%   ~/brain/2_projects/fao02/_dev_materials/release_notes/STYLE.md
%   ~/brain/2_projects/fao02/_dev_materials/prerelease_notes/STYLE.md
% British spelling. Title Case headings, sentence case bodies.
% booktabs rules only, no vertical rules, every table captioned and
% labelled.
% ============================================================

\documentclass[11pt,a4paper]{article}

% ------------------------------------------------------------
% LANGUAGE & ENCODING
% ------------------------------------------------------------
\usepackage[english]{babel}
\usepackage[utf8]{inputenc}
\usepackage[T1]{fontenc}        % REQUIRED: OT1 cannot set a, o, edh etc.

% ------------------------------------------------------------
% TYPOGRAPHY & SPACING
% ------------------------------------------------------------
\usepackage{microtype}           % Better kerning and justification
\usepackage[onehalfspacing]{setspace} % Line spacing
\usepackage{parskip}             % Paragraph spacing, no indent
\setlength{\footnotesep}{0.7\baselineskip}
\usepackage[hang,flushmargin]{footmisc} % Footnotes: no indent

% ------------------------------------------------------------
% PAGE LAYOUT
% ------------------------------------------------------------
\usepackage[a4paper,margin=1in]{geometry} % Clean margins
\usepackage{fancyhdr}
\pagestyle{fancy}
\fancyhf{} % clear all fields
% Head/foot fields are set after the delivery metadata is defined,
% further down. Inherited \lhead{Maase et al.} / \rhead{\today} removed
% deliberately: the date must not float on rebuild.
\renewcommand{\footrulewidth}{0.8pt}

% ------------------------------------------------------------
% GRAPHICS & TABLES
% ------------------------------------------------------------
\usepackage{graphicx}
\usepackage{booktabs}
\usepackage{tabularx}
\usepackage{xltabular}
\usepackage{array}
% Weighted full-width column: L{w} where the w's across a row sum to the column count.
\newcolumntype{L}[1]{>{\raggedright\arraybackslash\hsize=#1\hsize}X}
\usepackage{makecell}
\usepackage[section]{placeins} % Prevent floats before sections
\usepackage[table]{xcolor}

% ------------------------------------------------------------
% MATH & SYMBOLS
% ------------------------------------------------------------
\usepackage{amsmath,amssymb}

% ------------------------------------------------------------
% CITATION & QUOTING
% ------------------------------------------------------------
\usepackage{csquotes}
\usepackage{natbib}

% ------------------------------------------------------------
% HYPERLINKS
% ------------------------------------------------------------
\usepackage{hyperref}
\hypersetup{colorlinks=true, allcolors=black, urlcolor=[rgb]{0,0.25,0.55}}

% ------------------------------------------------------------
% NOTES
% ------------------------------------------------------------
\usepackage[colorinlistoftodos]{todonotes}

% ------------------------------------------------------------
% WATERMARK
%\usepackage{draftwatermark}
%\SetWatermarkText{DRAFT}
%\SetWatermarkScale{7}          % Size of the text (default: 1)
%\SetWatermarkLightness{0.9}    % 0 = black, 1 = invisible

% ------------------------------------------------------------
% DELIVERY METADATA --- fill all of these
% ------------------------------------------------------------
% IDENTIFIER. Composed from the facts below, never hashed:
%
%     <library>-<loa>-<scope>-<partition>-<YYYY-MM>[-<n>]
%
%   library    the model library that produced it: r2darts2, hydranet,
%              stepshifter, ...
%   loa        the VIEWS level of analysis: pgm (PRIO-GRID month) or
%              cm (country month). This carries BOTH the spatial unit
%              and the temporal resolution.
%   scope      the spatial extent: global, africa, ... Always present.
%              A delivery is not assumed global.
%   partition  calibration / validation / forecasting
%   YYYY-MM    month of delivery
%   n          sequence number, ONLY if a second delivery with the same
%              fields goes out in the same month
%
% CHECKSUM. \deliverychecksum is derived from every specification field
% below by ../checksum.py --- it is not chosen. Regenerate it whenever
% any of those fields changes:
%
%     python3 ../checksum.py . > /dev/null && python3 ../checksum.py .
%
% build.sh verifies it, so a note cannot quietly disagree with its own
% specification.
%
% \deliverydate is deliberately NOT \today. A note records what was sent
% on a given day; rebuilding it must not restamp it. Set it once and
% leave it.
%
% \deliverysha is the views-models commit THE DATA WAS PRODUCED AT ---
% not the commit of this document. Find it from the runs themselves (the
% MANIFEST beside each model's output records the sha it read its config
% at) and use the latest commit that touched any delivered model's
% configuration. It must be pushed, or the links below will not resolve.
%
% \cfg{model} links a model name to its hyperparameter configuration on
% GitHub, pinned to that commit.

\newcommand{\deliveryid}{r2darts2-pgm-global-calibration-2026-10}
\newcommand{\deliverychecksum}{60597598}
\newcommand{\deliverytitle}{Nine Deep-Learning Forecast Models}
\newcommand{\deliveryauthor}{Simon Polichinel von der Maase}
\newcommand{\deliverydate}{10-10-2026}
% --- SCOPE: what the data is --------------------------------
% Resolution and extent on each axis, stated in parallel. The level-of
% -analysis shorthand (pgm, cm) is not repeated here --- it is the second
% field of the identifier.
\newcommand{\deliverypartition}{Calibration}
\newcommand{\deliveryspatialres}{PRIO-GRID cell}
\newcommand{\deliveryspatialextent}{Global land --- 64,818 cells}
\newcommand{\deliverytemporalres}{Month}
\newcommand{\deliverytemporalextent}{Forecasts Jan 2018 to Dec 2021; trained on Jan 1990 to Dec 2017}
\newcommand{\deliverystructure}{13 rolling origins, each forecasting 36 months}
\newcommand{\deliverytargets}{State-based, non-state and one-sided fatalities}
% --- PROVENANCE: what produced it ---------------------------
\newcommand{\deliverylibrary}{views-r2darts2 0.2.4, wrapping Darts (\texttt{unit8co/darts})}
\newcommand{\deliveryframework}{views-pipeline-core 3.3.4}
\newcommand{\deliverybackend}{views-datafactory 1.13.0}
\newcommand{\deliverysha}{be7b7b6}
\newcommand{\deliverycompute}{Rented cloud GPUs (RunPod)}
% Link a model name to its hyperparameter config at the delivery commit.
\newcommand{\cfg}[1]{%
  \href{https://github.com/views-platform/views-models/blob/\deliverysha/models/#1/configs/config_hyperparameters.py}%
       {link}}

% ------------------------------------------------------------
% TITLE
% ------------------------------------------------------------
\usepackage{titling}

\pretitle{\vspace{2em}\centering\Huge\bfseries}
\posttitle{\par\vspace{1em}}
\preauthor{\centering\large}
\postauthor{\par\vspace{0.5em}}
\predate{\centering\small}
\postdate{\par\vspace{2em}}

\title{\deliverytitle}
\author{\deliveryauthor}
\date{\deliverydate}

% ------------------------------------------------------------
% RUNNING HEAD --- the identifier travels on every page, so a loose
% printed page can still be traced back to the delivery it came from.
% ------------------------------------------------------------
\lhead{\deliveryauthor}
\rhead{\deliverydate}
\lfoot{PRIO --- VIEWS \quad \texttt{\deliveryid}}
\rfoot{\thepage}

% ------------------------------------------------------------
% DOCUMENT STARTS
% ------------------------------------------------------------
\begin{document}

% ------------------------------------------------------------
% TITLE PAGE
% ------------------------------------------------------------
\begin{titlepage}
    \centering
    \vspace{3cm}

    {\Huge \textbf{VIEWS}\par}
    {\Huge Delivery Note \par}
    \vspace{0.6cm}
    {\LARGE \deliverytitle \par}
    \vspace{1.5cm}

    \noindent\makebox[\textwidth]{\Large\textemdash\quad$\diamond$\quad\textemdash}
    \vspace{1.5cm}

    \begin{tabular}{@{}r@{\quad}p{10cm}@{}}
        \textbf{Identifier}  & \texttt{\deliveryid} \\
        \textbf{Checksum}    & \texttt{\deliverychecksum} \\[0.35em]
        \textbf{Date}        & \deliverydate \\
        \textbf{Author}      & \deliveryauthor \\[0.9em]
        \multicolumn{2}{@{}l@{}}{\textbf{\textit{Scope}} \quad\small\textit{what the data is}} \\[0.2em]
        \textbf{Partition}            & \deliverypartition \\
        \textbf{Spatial resolution}   & \deliveryspatialres \\
        \textbf{Spatial extent}       & \deliveryspatialextent \\
        \textbf{Temporal resolution}  & \deliverytemporalres \\
        \textbf{Temporal extent}      & \deliverytemporalextent \\
        \textbf{Forecast structure}   & \deliverystructure \\
        \textbf{Targets}              & \deliverytargets \\[0.9em]
        \multicolumn{2}{@{}l@{}}{\textbf{\textit{Provenance}} \quad\small\textit{what produced it}} \\[0.2em]
        \textbf{Model library}        & \deliverylibrary \\
        \textbf{Framework}            & \deliveryframework \\
        \textbf{Data backend}         & \deliverybackend \\
        \textbf{Source commit}        & \texttt{views-models@\deliverysha} \\
        \textbf{Computed on}          & \deliverycompute \\
    \end{tabular}

    \vspace{0.8cm}
    \textit{PRIO -- Oslo}

    \vfill
    \includegraphics[width=0.25\textwidth]{figures/VIEWS_logo_tagline.png}
    \hspace{1cm}
    \includegraphics[width=0.25\textwidth]{figures/PRIO_Logo.png}

    \thispagestyle{empty}
\end{titlepage}

% ============================================================
% DOCUMENT BODY
% ------------------------------------------------------------
% There is NO fixed structure. Comment in what this delivery needs,
% comment out the rest, reorder freely. The headings below are
% suggestions drawn from past deliveries, not a required spine.
% ============================================================

\section{Context}
This Delivery Note accompanies \texttt{\deliveryid}, a set of conflict forecasts produced by nine
deep-learning models on the calibration partition at global PRIO-GRID month. It is the
authoritative description of what was handed over: the data itself carries no documentation, and
the folder is not self-explanatory.

The purpose of this note is to:
\begin{itemize}
    \item record what was delivered, in enough detail to identify it again later;
    \item describe the models that produced it, and what distinguishes them from one another;
    \item state how the files are organised and how they are intended to be consumed; and
    \item set out what is \emph{not} in the delivery, and what is not yet known about it.
\end{itemize}

\noindent This note covers:
\begin{itemize}
    \item the nine models, their architectures, and the published work they come from
          (Sections~\ref{sec:models}--\ref{sec:architectures});
    \item the objective functions they were trained on, including one written in-house
          (Section~\ref{sec:losses});
    \item the file layout, schema and units of the delivered forecasts
          (Section~\ref{sec:files});
    \item how the delivery is consumed, by the ensemble-construction tool or by hand
          (Section~\ref{sec:handover}); and
    \item what the delivery excludes, and the caveats that bear on using it
          (Sections~\ref{sec:excluded}--\ref{sec:caveats}).
\end{itemize}

\noindent This note explicitly does \emph{not} cover:
\begin{itemize}
    \item \textbf{any evaluation of the models.} Nothing here has been scored, ranked or compared.
          No metric has been computed beyond a check that the files are not corrupt;
    \item any recommendation as to which models to prefer, or whether to use them at all;
    \item the validation or forecasting partitions, or any live forecast; and
    \item the internal machinery that produced the run --- what it took to compute these
          forecasts is recorded separately, in
          \texttt{reports/postmortem\_darts\_calibration\_2026-10.md}.
\end{itemize}

\noindent The delivery is a fixed artefact. It is identified by the identifier and checksum on
the title page --- quote both --- and is not revised in place: any correction is issued as a new
delivery with a new identifier and a new checksum.


\section{The Models}\label{sec:models}
All nine are \textbf{off-the-shelf deep-learning models for time series}, taken from the Python
library \textbf{Darts} (\texttt{unit8co/darts}). The architectures are the library's, unmodified ---
each is a published model used as Darts implements it. What we supply is the data, the
configuration, and for six of the nine a loss function of our own (Section~\ref{sec:losses}).

They are \textit{global} models in the forecasting sense: a single set of parameters is pooled
across all 64,818 cells. The closest familiar analogue is a pooled panel estimator --- a cell with
almost no conflict history is still forecast using structure learned from the whole grid. Each
cell's own level and scale are restored separately, by the normalisation described in
Section~\ref{sec:architectures}.

\begingroup
\catcode`\_=12  % underscores are ordinary characters in this table
\begin{table}[h!]
\centering
\small
\renewcommand{\arraystretch}{1.2}
\caption{The nine delivered models.}
\label{tab:models}
\begin{tabular}{@{}p{2.9cm}p{1.7cm}p{1.9cm}p{2.7cm}p{3.1cm}p{1.5cm}@{}}
\toprule
\textbf{Model} & \textbf{Type} & \textbf{Loss} & \textbf{Predictors}
  & \textbf{Hyperparameters} & \textbf{Features} \\ \midrule
dark_river       & N-BEATS & MSE       & Outcomes only   & \cfg{dark_river}       & \qs{dark_river} \\
old_rules        & N-BEATS & Spotlight & Outcomes only   & \cfg{old_rules}        & \qs{old_rules} \\
blue_ocean       & N-BEATS & Spotlight & + 69 covariates & \cfg{blue_ocean}       & \qs{blue_ocean} \\ \addlinespace
silent_fox       & TSMixer & MSE       & Outcomes only   & \cfg{silent_fox}       & \qs{silent_fox} \\
dancing_monkey   & TSMixer & Spotlight & Outcomes only   & \cfg{dancing_monkey}   & \qs{dancing_monkey} \\
brave_heart      & TSMixer & Spotlight & + 69 covariates & \cfg{brave_heart}      & \qs{brave_heart} \\ \addlinespace
red_hawk         & TiDE    & MSE       & Outcomes only   & \cfg{red_hawk}         & \qs{red_hawk} \\
dark_necessities & TiDE    & Spotlight & Outcomes only   & \cfg{dark_necessities} & \qs{dark_necessities} \\
golden_eagle     & TiDE    & Spotlight & + 69 covariates & \cfg{golden_eagle}     & \qs{golden_eagle} \\
\bottomrule
\end{tabular}
\end{table}
\endgroup

\noindent The link columns open each model's configuration files on GitHub, fixed at the delivery
commit on the title page: \textbf{Hyperparameters} is \texttt{config\_hyperparameters.py}, which sets the
architecture, the loss and the training schedule; \textbf{Features} is
\texttt{config\_queryset.py}, which declares the predictors the model was given.

The names are arbitrary labels, assigned according to the \texttt{adjective\_noun} convention
used across views-models, and carry no meaning.

\textbf{Within each block of three, exactly two settings vary} --- the loss function, and whether
the model sees covariates. Everything else is identical: layer widths, depth, training epochs,
learning rate, early stopping. So inside a block, each row differs from the one above it in a
single setting. Across blocks the architecture changes and the comparison is a different one.

\begin{itemize}
    \item \textbf{Outcomes only} --- the model's inputs are the past values of the three outcomes
          and nothing else. Purely autoregressive. Six models: the first two rows of each block.
    \item \textbf{+ 69 covariates} --- the same three histories plus 69 predictors: ACLED event and
          fatality counts, land cover shares, population and built-up area, infant and child
          mortality, V-Dem governance indices, SHDI development indices, travel time to cities,
          mineral and petroleum deposits, growing season and harvested area, and terrain. Three
          models: the third row of each block.
\end{itemize}

\subsection*{How the Covariates Are Conditioned}

Three things happen to a covariate between the data store and the model. Only the last is ours.

\textbf{Spatial handling.} Nothing arrives natively on the PRIO-GRID. Population and built-up
surface are aggregated by summing $60\times60$ source raster pixels per cell; V-Dem and SHDI are
attached to cells by polygon-to-centroid spatial joins. The source documentation describes these as
opinionated decisions rather than mechanical transforms, and they are not reversible from the
delivered data.

\textbf{Temporal handling.} There is no feature engineering --- no lags, decay functions, ratios or
derived indices; the datafactory places cross-feature derivation in the model layer by design
(ADR-014 §6). But the series are not left as-is either. Annually observed quantities are extended
across their twelve months, by one of two per-source rules (ADR-014 §5): V-Dem and SHDI hold the
annual value flat; population and built-up surface interpolate linearly between five-yearly epochs.
Outside a source's observed range, values are \emph{manufactured}: zero before the first
observation, and the last observation held constant after it (ADR-052). Where a source's coverage
is narrower than the grid, the gap is filled with zeros rather than nulls, deliberately, because
consumers assume float arrays without missing-value handling (ADR-047). SHDI is the one exception:
cells its source never covered keep a null (ADR-042).

\textbf{In the model configuration.} All 72 columns --- the three outcomes and all 69 covariates
--- then pass through an \texttt{asinh} transform before training, declared in each model's
hyperparameter configuration (linked in Table~\ref{tab:models}). The transform is inverted for the
outcomes, which is why the delivered forecasts are in counts.

\subsection*{A Consequence Worth Knowing Before Using the Covariate Models}

The zero-fill rule above has a specific effect on this delivery, and it is not visible in the data:
\textbf{a manufactured zero is indistinguishable from an observed zero.}

The store publishes the real observation window per source. Checked against this delivery's
partition:

\begin{table}[h!]
\centering
\small
\renewcommand{\arraystretch}{1.2}
\caption{Source coverage against the calibration partition. Months outside a source's observed
window are manufactured, not measured.}
\label{tab:coverage}
\begin{tabular}{@{}p{3.4cm}p{3.6cm}p{3.6cm}p{3.4cm}@{}}
\toprule
\textbf{Source} & \textbf{Observed from} & \textbf{In training (1990--2017)} & \textbf{In test (2018--2021)} \\ \midrule
ACLED (7 covariates) & January 2020 & \textbf{none --- all zeros} & real from January 2020 \\
V-Dem                & January 1989 & real & real \\
SHDI                 & January 1989 & real & real \\
Population           & January 1989 & real & real \\
Built-up surface     & January 1989 & real & real \\
\bottomrule
\end{tabular}
\end{table}

\noindent \textbf{The seven ACLED covariates were identically zero throughout training.} They carry
no information for any month the three covariate-using models were fitted on, and then become
non-zero partway through the forecast window. A feature that is constant in training and varies in
test is not merely uninformative --- the model has learned that it is always zero.

This affects \texttt{blue\_ocean}, \texttt{brave\_heart} and \texttt{golden\_eagle} only. The six
outcomes-only models do not use covariates at all. The other four sources cover the whole partition.

\textbf{What this note does not address.} Whether any of this conditioning is appropriate --- the
fill rules, the zero-for-missing convention, the spatial joins, the \texttt{asinh} transform, or
the treatment of annual variables at monthly resolution --- is outside the scope of a delivery
note. Nothing here endorses those choices and nothing here has examined them. They are recorded so
that someone evaluating these models knows what the models were given. Two limitations of the
record itself are worth stating: provenance is held per \emph{source}, not per covariate, so
reconstructing what happened to any one of the 69 means assembling it from the source mapping, the
relevant ADR and the run ledger; and the population and built-up-surface series include projected
rather than measured epochs at recent dates, delivered identically to measured ones --- beyond this
partition's end, but worth knowing if the window is ever extended.

Note that the two splits are not the same split. Six models are outcomes-only and six use the
Spotlight loss, but they are different sixes: \texttt{blue\_ocean}, \texttt{brave\_heart} and
\texttt{golden\_eagle} have covariates \textit{and} Spotlight, while \texttt{dark\_river},
\texttt{silent\_fox} and \texttt{red\_hawk} have neither.

\textbf{There is no model combining MSE with the covariates.} Three configurations exist, not four.


\section{Model Architectures}\label{sec:architectures}
All three are feed-forward neural networks: information passes once from input to output, with no
internal memory carried between time steps. They differ from one another in how the layers are
arranged, described below.

\subsection*{N-BEATS}

A deep stack of fully connected layers with backward and forward residual links. It builds the
forecast as a sum of learned basis functions, which is what makes it partly interpretable. Designed
for univariate point forecasting and benchmarked on the M3, M4 and TOURISM competition datasets.

\noindent Implemented in Darts as \href{https://unit8co.github.io/darts/generated_api/darts.models.forecasting.nbeats.html}{\texttt{NBEATSModel}};
see \citet{oreshkin2020nbeats}.


\subsection*{TSMixer}

Dense layers only, applied alternately along the time axis and across the predictor variables, so
the model mixes temporal structure and cross-variable structure in turn. Built for the multivariate
case.

\noindent Implemented in Darts as \href{https://unit8co.github.io/darts/generated_api/darts.models.forecasting.tsmixer_model.html}{\texttt{TSMixerModel}};
see \citet{chen2023tsmixer}.


\subsection*{TiDE}

An encoder--decoder, again dense layers only. It compresses the history and the predictors into a
single code, then decodes the whole 36-month path from it in one pass.

\noindent Implemented in Darts as \href{https://unit8co.github.io/darts/generated_api/darts.models.forecasting.tide_model.html}{\texttt{TiDEModel}};
see \citet{das2023tide}.



\section{Loss Functions}\label{sec:losses}

\textbf{MSE} --- ordinary squared error, every cell-month weighted equally. The library's own
implementation, a thin wrapper around PyTorch's. Used by \texttt{dark\_river}, \texttt{silent\_fox}
and \texttt{red\_hawk}.

\textbf{Spotlight} (\texttt{SpotlightLossLogcosh}) --- in-house, not off the shelf. A three-part
objective computed after an \texttt{asinh} transform. It scores the \textit{shape} of the error only
in cells where the observation or the forecast exceeds one fatality; separately penalises the
\textit{level} in cells whose observation exceeds one fatality; and separately penalises
forecasting conflict in cells that stayed below it. It uses \texttt{log cosh} rather than squared
error, so a single very large observation does not dominate the gradient. Used by the other six.

\subsection*{Development and Validation Status}

Spotlight is not a one-off function but the production member of a sustained line of work. The
first variant was committed in March 2026 and the one used here at the end of April; development
continued into late August. The family now comprises six Spotlight variants and three related
objectives --- roughly 2,250 lines across nine modules --- attracting some 216 commits over six
months. The library's own documentation sets out five earlier generations of the design, recording
for each what it addressed and what it exposed. By any ordinary measure this represents a
considerable investment of development effort.

The objective takes a single parameter, and it is derived rather than fitted: the threshold
separating event from quiet cells is $\mathrm{asinh}(1) \approx 0.88$, exactly one fatality on the
transformed scale.

What the record does not contain is any formal account of it. No citation appears in the source,
its documentation, the commit history or the architecture decision records; two other objectives in
the same library do cite published work, so the omission is particular to this family rather than a
habit of the codebase. Neither is it covered by tests. The library's suite exercises eight of its
objectives and this is not among them, so nothing checks its gradients, compares it against a
reference implementation, or establishes that the code computes what the written account describes
--- this for the objective on which 37 of the platform's 42 darts models are trained.

Six months of concentrated development across five successive generations is not the profile of an
untried design, and the written account of its construction is detailed. It is reasonable to infer
that the objective has been exercised extensively, if informally, during that period. Informal
exercise is not validation, however, and the distinction is what this note is obliged to record:
what is absent is not the work but the evidence of it. A derivation a third party can follow, and
tests that fail when the implementation is wrong, remain outstanding and should be produced for an
objective in production use at this scale.

On that basis this note takes no position on whether the objective is sound. It may be excellent,
adequate, or quietly wrong, and nothing in the available record distinguishes those possibilities.
Six of the nine delivered models were trained on it and three on ordinary squared error, which
bears on how much weight a comparison between those two groups can carry.


\section{The Delivered Files}\label{sec:files}
\textbf{One folder per model.} Inside each: thirteen Parquet files --- a columnar table format
that opens in one call from both Python (\texttt{pandas.read\_parquet}) and R
(\texttt{arrow::read\_parquet}) --- named \texttt{\dots\_00} through
\texttt{\dots\_12}, and a short text file recording when the model ran, which version of the
software produced it, and which configuration it used.

\textbf{The thirteen files are thirteen forecast origins}, stepping forward one month at a time. An
origin is the last month the model was allowed to see; from there it forecasts the next 36 months.

\begin{table}[h!]
\centering
\renewcommand{\arraystretch}{1.2}
\caption{The thirteen forecast origins.}
\label{tab:origins}
\begin{tabular}{@{}p{3cm}p{4cm}p{5cm}@{}}
\toprule
\textbf{File} & \textbf{Last month seen} & \textbf{Months forecast} \\ \midrule
\texttt{\dots\_00} & December 2017 & January 2018 -- December 2020 \\
\texttt{\dots\_01} & January 2018  & February 2018 -- January 2021 \\
\multicolumn{3}{@{}c@{}}{\dots} \\
\texttt{\dots\_12} & December 2018 & January 2019 -- December 2021 \\
\bottomrule
\end{tabular}
\end{table}

Within every file the earliest month is horizon 1 and the latest horizon 36, so a given target
month sits at a different horizon in each file containing it. Metrics computed per origin keep
those apart; metrics pooled across origins mix them.

\textbf{Each file holds 2,333,448 rows} --- 36 months $\times$ 64,818 cells --- and five columns.

\begin{table}[h!]
\centering
\renewcommand{\arraystretch}{1.2}
\caption{Columns in every file.}
\label{tab:columns}
\begin{tabular}{@{}p{4.5cm}p{9.5cm}@{}}
\toprule
\textbf{Column} & \textbf{Contents} \\ \midrule
\texttt{month\_id}        & Month \\
\texttt{priogrid\_id}     & Grid cell \\
\texttt{pred\_lr\_ged\_sb} & Forecast state-based fatalities \\
\texttt{pred\_lr\_ged\_ns} & Forecast non-state fatalities \\
\texttt{pred\_lr\_ged\_os} & Forecast one-sided fatalities \\
\bottomrule
\end{tabular}
\end{table}

\textbf{\texttt{priogrid\_id} is a cell, not a coordinate.} PRIO-GRID is a
$720 \times 360$ grid of half-degree cells covering the globe, numbered from 1 in the
south-west corner and running east before north. To place a row on a map:

\begin{table}[h!]
\centering
\renewcommand{\arraystretch}{1.3}
\caption{Converting \texttt{priogrid\_id} to the centre of its cell.}
\label{tab:priogrid}
\begin{tabular}{@{}p{4.6cm}p{9.4cm}@{}}
\toprule
\textbf{Quantity} & \textbf{From \texttt{priogrid\_id}} \\ \midrule
Row (south to north) & $\texttt{row} = \lfloor (\texttt{priogrid\_id} - 1) / 720 \rfloor$ \\
Column (west to east) & $\texttt{col} = (\texttt{priogrid\_id} - 1) \bmod 720$ \\
Latitude  & $-90 + (\texttt{row} + 0.5) \times 0.5$ \\
Longitude & $-180 + (\texttt{col} + 0.5) \times 0.5$ \\
\bottomrule
\end{tabular}
\end{table}

\noindent In R, for a data frame \texttt{d}:

\begin{quote}
\ttfamily\small
d\$lat <- -90  + (((d\$priogrid\_id - 1) \%/\% 720) + 0.5) * 0.5 \\
d\$lon <- -180 + (((d\$priogrid\_id - 1) \%\%  720) + 0.5) * 0.5
\end{quote}

\noindent The coordinates are cell \emph{centres}; each cell spans $0.5^\circ$ in both
directions. As a check, \texttt{priogrid\_id} 148759 --- the cell discussed in
Section~\ref{sec:caveats} --- resolves to $13.25^\circ$N, $39.25^\circ$E, in Tigray.
The 64,818 cells delivered here are the land subset of the 259,200 in the full grid.

\textbf{Units are fatality counts.} A value of 2.4 means 2.4 fatalities; values are continuous
rather than integer. The transformations applied during fitting are intended to be inverted
before these files are written, so no back-transformation should be applied to them;
Section~\ref{sec:transforms} states what is and is not established about that.

\textbf{Observed values are not included, and are not needed from here.} Evaluation joins these
forecasts to observed fatalities on \texttt{month\_id} and \texttt{priogrid\_id}; the
ensemble tool retrieves those itself, and they are otherwise available from the
datafactory. The file naming used here already matches what that tool expects.

Nine models $\times$ thirteen origins = \textbf{117 files, about 2\,GB} in total.


\section{Consuming the Delivery}\label{sec:handover}
\subsection*{Observed Values}

The evaluation joins forecasts to observed fatalities on \texttt{month\_id} and
\texttt{priogrid\_id}. Those observations are not in this folder and do not need to be: the
ensemble-construction tool retrieves them itself, and they are in any case available from the
datafactory. The file naming here already matches what that tool expects.

\subsection*{Computing Metrics Directly}

The instruction is: \textit{iterate over the thirteen files per model, join each to the observed
values on} \texttt{month\_id} \textit{and} \texttt{priogrid\_id}\textit{, and compute point metrics
per origin so that horizon effects stay separable.} Parquet reads in one call ---
\texttt{pandas.read\_parquet} in Python, \texttt{arrow::read\_parquet} in R.

Two things to tell them explicitly, because both will otherwise be guessed wrong:

\begin{enumerate}
    \item forecasts are counts, not logs --- no back-transformation;
    \item there is no distribution --- no distributional scores.
\end{enumerate}


\section{Not Included in This Delivery}\label{sec:excluded}
\begin{table}[h!]
\centering
\renewcommand{\arraystretch}{1.3}
\caption{Not included in this delivery.}
\label{tab:not-included}
\begin{tabular}{@{}p{4cm}p{10cm}@{}}
\toprule
\textbf{Item} & \textbf{Where it is instead, or why not} \\ \midrule
Observed values     & Fetched by the ensemble tool, or from the datafactory
                      (Section~\ref{sec:files}) \\
Evaluation metrics  & None computed --- see Section~\ref{sec:caveats} \\
Training logs       & Not shipped. \texttt{calibration\_log.txt}, which the ensemble tool can
                      optionally read to resolve run timestamps, is among them; timestamps can be
                      passed explicitly instead \\
Other partitions    & Calibration only --- no validation partition, no true forecasts \\
Model artefacts     & No weights, no checkpoints \\
\bottomrule
\end{tabular}
\end{table}


\section{Caveats and Notes}\label{sec:caveats}

\textbf{No model here has a predictive distribution.} These are point forecasts --- one number per
cell-month, no draws, no quantiles, no intervals. All nine are configured as deterministic
predictors, so a different export would not recover a distribution; it was never produced. Nothing
that scores a \textit{single model's} uncertainty --- CRPS of its predictive density, interval
coverage, density calibration --- can be computed from these files.

That is a limitation of this delivery against the direction the VIEWS Prediction Challenge sets
out, which argues that for a field this skewed a distribution is closer to what is needed than a
point estimate. These nine do not provide one.

\begin{quote}
\small
Hegre, H., Vesco, P., Colaresi, M. et al. (2025), \textit{The 2023/24 VIEWS Prediction Challenge:
Predicting the number of fatalities in armed conflict, with uncertainty}, Journal of Peace
Research. \texttt{doi:10.1177/00223433241300862}
\end{quote}

One qualification, since it would otherwise look like a contradiction in the output: the ensemble
tool can still produce a CRPS. It does so by treating the member models' point forecasts as an
empirical sample and scoring the spread across models. That is a property of the ensemble, not of
any model in it, and it is not the same quantity as a CRPS computed from a model's own posterior.
Worth knowing before the two are compared.

\textbf{Nothing here has been evaluated.} I checked only that the files are not broken: no
negatives, no missing values, no duplicate cell-months, the expected row counts, and totals within
the same order of magnitude as observed. That rules out a failed back-transformation or a collapsed
model and establishes nothing else.

\textbf{It is nine models, not eleven.} Two others were set up to draw 100 samples per forecast,
making them probabilistic counterparts of \texttt{dancing\_monkey} and \texttt{dark\_necessities}.
A point delivery needs one value per cell, and at one sample their configurations are
byte-identical to those siblings --- the same model twice. Running them as intended is blocked by
memory: 100 samples for 64,818 cells across thirteen origins needs far more RAM than any machine we
have.

\textbf{One observation will dominate any squared-error metric.} December 2021, cell 148759 has
\textbf{113,395} observed fatalities; the next largest anywhere in the window is 5,395. That month
falls in only one of the thirteen origins (\texttt{\dots\_12}), where that single cell is about
41\% of observed fatalities. I do not know whether it is a real event or a data error. Either way
it is better settled before the numbers are produced than explained afterwards.

\textbf{The software is pinned; the data is not.} The title page records the three versions that
determine the output --- the model library, the framework and the data backend --- along with the
\texttt{views-models} commit the configurations were read at. Between them those fix the code.
They do not fix the \textit{input}: the data backend served this run from a live store with no
snapshot identifier, and the provenance record written beside the fetched data shows drift
detection did not run. So re-running these nine models at these versions would reproduce the
method, and would not guarantee reproducing the numbers, because the underlying observations may
have been revised in the meantime. Worth knowing before anyone treats a re-run as a check.

\textbf{Column naming.} These files use the canonical target names,
\texttt{pred\_lr\_ged\_sb} and its two siblings. The September delivery used
\texttt{pred\_lr\_sb\_best} instead, so the two are not directly comparable without a rename.
Renaming either set takes minutes and needs no re-run.

\textbf{Comparison with the September delivery.} You were sent three variants per model then --- a
mean, a 95th percentile and a conditional mean --- each computed from that model family's
predictive distribution. These nine have no distribution, so only the mean has an equivalent here.
If the evaluation work has been built on one of the other two, the two deliveries are not
like-for-like. A correction while I am here: my September note said the full distributions had been
retained so alternative summaries could be produced without re-running. They were not kept.


% ------------------------------------------------------------
% BIBLIOGRAPHY --- uncomment ONLY if the note actually cites
% something, otherwise it renders an empty References section.
% Add a bibtex pass to build.sh if you enable this.
% ------------------------------------------------------------
% \bibliographystyle{apalike}
% \bibliography{conflict}

\end{document}
